%% title: Python Implementation
%% category: misc

\begin{minted}{python}
### MTSM implementation in Python

tier_list = ['person', 'number', 'gender'] 

# A syntactic node. 
class snode:
    def __init__(self, left=None, right=None, mother=None):
        self.vals = {} # list of tier values
        for tier in tier_list:
            self.vals[tier] = None
        self.left = left # left daughter
        self.right = right # right daughter
        self.mother = mother # mother
    def __repr__(self):
        return f"snode(vals={self.vals!r})"

# A transducer which gives values to non-terminals, given a transition function. 
def run_mtsm(input_stream, tau):
    output_stream = []
    right = None; # assumes the bottommost right node is a lexical root unvalued for any features.
    for left in reversed(input_stream):
        mother = tau(left, right) # tau is the transition function
        output_stream.insert(0, right)
        output_stream.insert(0, left)
        right = mother
    output_stream.insert(0, right)
    return output_stream

    # Input: list of left nodes from top to bottom
    # Output: list of nodes in "right[left right[left right[...]]]" order (all right daughters are even-indexed). 

    # Important: treats specifiers as terminals, and assumes this process has already run on them. 
    # Important: assumes the strict LCA, and operates over constituency trees, not dependency trees.

# A helper to retrieve the value of a given tier on a given node. 
def val(node, tier):
    return None if node is None else node.vals[tier]

# A helper to check that a tree is well-formed. 
def validate_tree(tree):
    if len(tree) % 2 == 0: raise ValueError('tree has even number of nodes')
    if any(val(tree[-1], t) is not None for t in tier_list):
        raise ValueError('bottom right node is not empty')

# An evaluator which checks if a given tree could have been created with an MTSM or not, given a specific tau. 
def eval_tauful(tree, tau):
    validate_tree(tree)
    for i in range(0, len(tree)-1, 2):
        mother, left, right = tree[i], tree[i+1], tree[i+2]
        if any(tau(left, right).vals[t] != val(mother, t) for t in tier_list):
            return False
    return True

# An evaluator which checks if there exists a possible tau that could have created this MTSM. 
def eval_tauless(tree):
    validate_tree(tree)
    seen = {} # dict of tier, leftval, rightval -> motherval, to enforce independence
    for i in range(0, len(tree)-1, 2):
        mother, left, right = tree[i], tree[i+1], tree[i+2]
        for tier in tier_list:
            l, r, m = val(left, tier), val(right, tier), val(mother, tier)
            if m != l and m != r: 
                return False # check conservation
            if seen.setdefault((tier, l, r), m) != m: 
                return False # check independence
    return True

# An example tau, which first checks a person feature hierarchy, and then defaults to the left node's feature unless it's unvalued. (This tau is associative.)
person_hierarchy = {'1': 0, '2': 1, '3': 2}
def tau(left, right):
    mother = snode(left, right)
    for t in tier_list:
        l, r = val(left, t), val(right, t)
        if l is None or r is None:
            mother.vals[t] = r if l is None else l
        elif t == 'person' and person_hierarchy[r] < person_hierarchy[l]:
            mother.vals[t] = r
        else: 
            mother.vals[t] = l
    return mother

# A given tree with values is generatable by an MTSM with a conservative and independent transition function tau iff there exists a tree and an initialization of left-node values which produces the output stream via tau. 

#############
# Other notes:
    # The current MTSM mandates the following properties of the transition function tau:
        # 1. conservation: for each tier, the mother value is one of the daughter values.
        # 2. independence: for each ordered pair of daughter values, there is only one possible mother value.

    # Other properties which the current MTSM doesn't mandate but which might be true:
        # 3. associativity: tau(tau(left, right), third) = tau(left, tau(right, third)).
        # 4. propensity for valuedness: if at least one daughter is valued for a feature, the mother cannot be unvalued for that feature. 

    # Dependency trees kind of assume by default that head features overpower specifier features, whereas constituency trees kind of make the opposite assumption, if you assume that everything overtakes the last thing. 

    # Hierarchies are probably universal, but which parts of them are actualized might not be -- e.g. maybe in some languages, 2=3 and 1>both, but in other languages 1=2 and both>3.

    # The current MTSM is more restrictive than OTSL for two reasons, both underspecified in OTSL:
        # 1. The OTSL tier projection function can potentially condition on any of the current node's features. For example, in TSL syntax, Kenneth often uses this as a tool to project determiners to an agreement tier depending on their case or EPP features. 
        # 2. The OTSL tier projection function can potentially condition on any of the last-projected-node's features (it's unclear). 

        # Crucially, this is because of the MTSM's property of Independence, which prevents different tiers from influencing each other. 

        # Without Independence, the MTSM could be equivalent to OTSL. 

        # However, it's important to note that this would fix (2) but not fully (1). OTSL and the MTSM can both condition on nodes below them, and this is fine, because those nodes' feature values have already been decided (i.e., do not condition on anything above them). But once you allow for conditioning on different features of your own node, you need to enforce a strict ordering *between* tiers/feature types. 

        # Otherwise, you will have both nodes that condition on features below them, AND nodes that condition on features above them, and this could get quite complex, and doesn't seem like something language would do. 
\end{minted}
